\scanpage{382}
\mspage{121}
% archive sheet number: 381

Bien entendu, les flèches
\[
C_\infty(\underline t')\mathrel{\unsure{-|\to}}C_\infty(\underline t)
\tag*{$(*)$}
\]
dans $\mathrm{Hot}$, associées aux flèches $\underline t\mathrel{-}\underline t'$ dans $\underline{\mathrm{Hom\,cof}}(\mathbb N^\circ,T)$, sont en fait des flèches dans $(\mathrm{Cat})$ elle-même -- et $=$ des foncteurs pleinement fidèles, dans le cas ``géométrique'' et pour les deux flèches canoniques
\diagram{\begin{tikzcd}[ampersand replacement=\&,row sep=small,column sep=small]
\&\underline t*\underline t'\arrow[dl]\arrow[dr,dash]\&\\
\underline t\&\&\underline t'\,.
\end{tikzcd}}

Elles sont isomorphiques \uline{dans $\mathrm{Hot}$}, parce qu'elles \deleted{\unsure{limites inductives}} définissent un isomorphisme sur les ind-objets associés de $\mathrm{Cat\,hot}$, donc aussi pour les ind-objets associés de $\mathrm{Hot}$ et \deleted{de $\mathrm{Hot}$} dans $\mathrm{Hot}_0^{\unsure{\wedge|1}}$, d'où résulte que
\[
\mathrm{Hom}_{\mathrm{Hot}}(Y,C_\infty(\underline t'))\mathrel{-}
\mathrm{Hom}_{\mathrm{Hot}}(Y,C_\infty(\underline t))
\]
associé à $(*)$ est une \deleted{\unsure{isom.}} \inserted{bij.}, \uline{si $Y$ est finie}$^{*}$, cette flèche s'identifiant alors à
% slanted
\margin{Ce qui suffit à assurer que $(*)$ est iso dans $\mathrm{Hot}$\unsure{, ;| ;} $\mathrm{Hot}\to\mathrm{Hot}_0^\wedge$ est \uline{conservatif} !}
\[
\underset{n}{\underrightarrow{\lim}}\,\mathrm{Hom}_{\mathrm{Hot}}(Y,C_{t'_n}\deleted{\unsure{(t'_n)}})
\longrightarrow\underset{n}{\underrightarrow{\lim}}\,\mathrm{Hom}_{\mathrm{Hot}}(Y,C_{t_n}\deleted{\unsure{(t_n)}}).
\]
(Car \inserted{$Z\mapsto$}$\mathrm{Hom}_{\mathrm{Hot}}(Y,Z)$ commute aux limites inductives filtrantes, si $Y$ est une catégorie \uline{finie}. \deleted{\unsure{)}} Cela provient du fait qu'une $\underrightarrow{\lim}$ filtrante de catégories de \unsure{Thomason-Kan|Thomason Kan} est itou\ldots)

Mais j'ignore par contre si $(*)$ est iso dans $\mathrm{Cat\,hot}$. \unsure{Sachant|Seulement} que pour une catégorie \uline{finie} $Y$, elle induit
\[
\mathrm{Hom}_{\mathrm{Cat\,hot}}(Y,C_\infty(\underline t'))\simeq
\mathrm{Hom}_{\mathrm{Cat\,hot}}(Y,C_\infty(\underline t)),
\]
les deux membres s'identifiant alors à
